\documentclass[../main.tex]{subfiles}
\begin{document}

% ============================================================
% TikZ: Δίτροχο μοντέλο οχήματος -- συστήματα συντεταγμένων
%        Slip angle ορισμός -- Μεταφορά φορτίου
%
% Χρήση: \subfile{tikz/VehicleModelDiagram} μέσα σε figure env.
%        Ή compile standalone με:
%   \documentclass[tikz,border=5pt]{standalone}
% ============================================================

% ---------------------------------------------------------------
% Σχήμα 1: Δίτροχο μοντέλο -- κάτοψη (Top View)
% ---------------------------------------------------------------
\begin{figure}[htbp]
\centering
\begin{tikzpicture}[scale=1.0,
    wheel/.style={rectangle, draw=black, fill=gray!30, minimum width=0.25cm, minimum height=0.7cm, rounded corners=1pt},
    force/.style={->, thick, blue!70!black},
    frameaxis/.style={->, thick},
    carframe/.style={->, thick, red!70!black},
    tyreframe/.style={->, thick, green!50!black},
    dim/.style={<->, thin, gray},
    label/.style={font=\small}
]

%% --- Karosserie (car body) ---
\draw[fill=gray!15, draw=black, rounded corners=3pt]
    (-1.0, -0.55) rectangle (1.0, 0.55);
\node[font=\scriptsize] at (0,0) {Σώμα};

%% --- Τροχοί (wheels) ---
% FL (Front Left)
\node[wheel] (FL) at (-0.65,  0.85) {};
\node[font=\scriptsize, above=1pt] at (FL.north) {FL};

% FR (Front Right)
\node[wheel] (FR) at (-0.65, -0.85) {};
\node[font=\scriptsize, below=1pt] at (FR.south) {FR};

% RL (Rear Left)
\node[wheel] (RL) at ( 0.65,  0.85) {};
\node[font=\scriptsize, above=1pt] at (RL.north) {RL};

% RR (Rear Right)
\node[wheel] (RR) at ( 0.65, -0.85) {};
\node[font=\scriptsize, below=1pt] at (RR.south) {RR};

%% --- Αδρανειακό σύστημα (Global frame) ---
\begin{scope}[shift={(-2.5, -1.8)}]
    \draw[frameaxis, black] (0,0) -- (0.8,0) node[right, font=\small] {$X$ (forward)};
    \draw[frameaxis, black] (0,0) -- (0,0.8) node[above, font=\small] {$Y$ (left)};
    \node[font=\scriptsize, below left] at (0,0) {Global};
\end{scope}

%% --- Σύστημα σώματος (Car frame) -- ευθυγραμμισμένο με τον άξονα οχήματος ---
%% Το σώμα κινείται προς τα αριστερά στο παράδειγμα (β>0)
\begin{scope}[shift={(0,0)}, rotate=10]
    \draw[carframe] (0,0) -- (0.7,0) node[right, font=\scriptsize, text=red!70!black] {$x$};
    \draw[carframe] (0,0) -- (0,0.7) node[above, font=\scriptsize, text=red!70!black] {$y$};
\end{scope}
\node[font=\scriptsize, text=red!70!black] at (0.6, 0.5) {Car frame};

%% --- β angle arc ---
\draw[->, thin, purple] (0.45,0) arc[start angle=0, end angle=10, radius=0.45];
\node[font=\scriptsize, purple] at (0.7, 0.15) {$\beta$};

%% --- Γωνία δ (steer angle) -- εμπρός τροχός FL ---
\begin{scope}[shift={(-0.65, 0.85)}, rotate=20]
    \draw[tyreframe, green!60!black] (0,0) -- (0.35,0);
\end{scope}
\draw[->, thin, green!40!black] (-0.65+0.25, 0.85) arc[start angle=0, end angle=20, radius=0.25];
\node[font=\scriptsize, green!40!black] at (-0.2, 1.1) {$\delta$};

%% --- Velocity vector ---
\draw[->, ultra thick, orange!80!black] (0,0) -- (-0.6, 0.25)
    node[left, font=\scriptsize, orange!80!black] {$\mathbf{v}$};

%% --- Wheelbase & Track labels ---
\draw[dim] (-0.65, 1.15) -- (0.65, 1.15);
\node[font=\scriptsize, gray, above] at (0, 1.15) {$L$ (wheelbase)};
\draw[dim] (-1.05, 0.85) -- (-1.05, -0.85);
\node[font=\scriptsize, gray, left] at (-1.05, 0) {$t$ (track)};

%% --- Center of Mass ---
\filldraw[black] (0,0) circle (2pt);
\node[font=\scriptsize, below right] at (0.05,-0.05) {CG};

\end{tikzpicture}
\caption{Κάτοψη δίτροχου μοντέλου. Φαίνονται οι τέσσερις τροχοί, το αδρανειακό σύστημα (μαύρο), το σύστημα σώματος (κόκκινο) υπό γωνία πλαϊνής ολίσθησης $\beta$, και η γωνία διεύθυνσης εμπρός τροχού $\delta$ (πράσινο).}
\label{fig:twotrack_topview}
\end{figure}


% ---------------------------------------------------------------
% Σχήμα 2: Ορισμός γωνίας ολίσθησης (Slip Angle α)
% ---------------------------------------------------------------
\begin{figure}[htbp]
\centering
\begin{tikzpicture}[scale=1.2,
    wheel/.style={rectangle, draw=black, fill=gray!30, minimum width=0.3cm, minimum height=1.0cm, rounded corners=2pt},
    vec/.style={->, thick},
    annot/.style={font=\small}
]

%% Τροχός
\node[wheel] (W) at (0,0) {};

%% Άξονας κύλισης τροχού (heading direction)
\draw[vec, black!70] (0,-0.8) -- (0, 1.2) node[above, annot] {κατεύθυνση τροχού};

%% Πραγματικό διάνυσμα ταχύτητας (υπό γωνία α)
\draw[vec, blue!80!black, very thick] (0,0) -- (0.6, 1.0)
    node[right, annot, blue!80!black] {$\mathbf{v}_{\text{wheel}}$};

%% Γωνία α
\draw[->, thin, red] (0, 0.5) arc[start angle=90, end angle=59, radius=0.5];
\node[annot, red] at (0.28, 0.68) {$\alpha$};

%% Vy component
\draw[dashed, gray] (0, 1.0) -- (0.6, 1.0);
\draw[vec, gray!60] (0, 0) -- (0.6, 0) node[below, annot, gray!60] {$v_y$};
\draw[vec, gray!60] (0.6, 0) -- (0.6, 1.0) node[right, annot, gray!60] {$v_x$};

\node[annot, below=8pt] at (0,0) {$\alpha = \arctan\!\left(\dfrac{v_y}{v_x}\right)$};

\end{tikzpicture}
\caption{Ορισμός γωνίας ολίσθησης $\alpha$. Η διαφορά μεταξύ της πραγματικής ταχύτητας $\mathbf{v}_{\text{wheel}}$ και της κατεύθυνσης κύλισης του τροχού δίνει τη γωνία $\alpha$, που αντιστοιχεί σε πλευρική δύναμη $F_y$.}
\label{fig:slip_angle}
\end{figure}


% ---------------------------------------------------------------
% Σχήμα 3: Μεταφορά φορτίου (Weight Transfer)
% ---------------------------------------------------------------
\begin{figure}[htbp]
\centering
\begin{tikzpicture}[scale=0.9,
    spring/.style={decorate, decoration={coil, amplitude=4pt, segment length=5pt}},
    vec/.style={->, very thick},
    dim/.style={<->, thin, gray}
]

%% Έδαφος
\draw[thick, gray] (-3.2, 0) -- (3.2, 0);
\fill[gray!20] (-3.2,-0.15) rectangle (3.2, 0);

%% Τροχοί (πλευρική άποψη -- 2 τροχοί για longitudinal transfer)
\filldraw[gray!40, draw=black] (-2.0, 0.25) circle (0.25);
\filldraw[gray!40, draw=black] ( 2.0, 0.25) circle (0.25);

%% Σώμα (πλευρική άποψη)
\draw[fill=gray!15, draw=black, rounded corners=3pt]
    (-2.2, 0.5) rectangle (2.2, 1.1);

%% Κέντρο μάζας
\filldraw[red] (0, 1.1) circle (3pt);
\node[font=\small, above right, red] at (0.1, 1.1) {CG};
\draw[dashed, gray] (0, 0.5) -- (0, 1.4);

%% Ύψος CG
\draw[dim] (-2.5, 0.5) -- (-2.5, 1.1);
\node[gray, font=\small, left] at (-2.5, 0.8) {$h_{\text{CG}}$};

%% Επιτάχυνση (ax) -- μπροστά
\draw[vec, orange!90!black] (0, 1.6) -- (1.5, 1.6)
    node[right, font=\small] {$a_x > 0$ (επιτάχυνση)};

%% Φορτία στατικά
\draw[vec, blue!70] (-2.0, 0.5) -- (-2.0, -0.7) node[below, font=\small, blue!70] {$F_{z,\text{front}}$};
\draw[vec, blue!70] ( 2.0, 0.5) -- ( 2.0, -0.9) node[below, font=\small, blue!70] {$F_{z,\text{rear}}$};

%% Μεταβολή φορτίου (ΔFz)
\draw[vec, red!70, dashed] (-2.0, -0.5) -- (-2.0, -0.1);
\node[font=\scriptsize, red!70, left] at (-2.0, -0.3) {$-\Delta F_z$};
\draw[vec, red!70, dashed] (2.0, -0.5) -- (2.0, -0.9);
\node[font=\scriptsize, red!70, right] at (2.0, -0.7) {$+\Delta F_z$};

%% Wheelbase
\draw[dim] (-2.0, -1.2) -- (2.0, -1.2);
\node[gray, font=\small, below] at (0, -1.2) {$L$};

%% Εξίσωση
\node[font=\small] at (0, -1.7) {$\Delta F_{z} = \dfrac{m \cdot a_x \cdot h_{\text{CG}}}{L}$};

\end{tikzpicture}
\caption{Διαμήκης μεταφορά φορτίου κατά την επιτάχυνση ($a_x > 0$). Το φορτίο μεταφέρεται από τον εμπρός άξονα (−$\Delta F_z$) στον πίσω (+$\Delta F_z$). Αντίστοιχα ισχύει για πλευρική μεταφορά υπό $a_y$.}
\label{fig:weight_transfer}
\end{figure}

\end{document}
